% Luc Destombe --- licence CC BY-NC-SA 4.0
% https://creativecommons.org/licenses/by-nc-sa/4.0/deed.fr
% Fichier autonome : compiler avec pdflatex, deux fois (signets, nombre de pages).
\documentclass[10pt,a4paper,twocolumn]{article}
\makeatletter
\newif\iflycee@unecolonne
\newif\iflycee@palatino
\lycee@palatinotrue

\RequirePackage[utf8]{inputenc}
\RequirePackage[T1]{fontenc}
\RequirePackage[french]{babel}
\RequirePackage{etoolbox}

\newcommand{\ShorthandsOff}{\@ifundefined{shorthandoff}{}{\shorthandoff{;:!?}}}
\newcommand{\ShorthandsOn}{\@ifundefined{shorthandon}{}{\shorthandon{;:!?}}}
\ShorthandsOff
\AtBeginDocument{\ShorthandsOn}
\AtBeginEnvironment{tikzpicture}{\ShorthandsOff}

\RequirePackage{geometry}
\RequirePackage{fancyhdr}
\RequirePackage{lastpage}
\RequirePackage{multicol}
\RequirePackage{array}
\RequirePackage{enumitem}
\RequirePackage{xcolor}
\RequirePackage{needspace}

\RequirePackage{amsmath,amssymb}
\RequirePackage{mathpazo}
\DeclareMathSizes{10.95}{10.95}{8}{5.5}
\begingroup\catcode`\_=\active \catcode`\^=\active
\gdef_#1{\sb{\scriptscriptstyle #1}}
\gdef^#1{\sp{\scriptscriptstyle\lycee@moinsepais #1}}
\endgroup
\newcommand\lycee@moins{\mathbin{%
  \setbox\z@\hbox{$\m@th\scriptscriptstyle\lycee@moinsorig$}%
  \dimen@=.55\wd\z@ \advance\dimen@-.5pt
  \hbox to.75\wd\z@{\hss$\m@th\scriptscriptstyle\vcenter{\hbox{%
    \tikz\draw[line width=.5pt,line cap=round](0,0)--(\the\dimen@,0);}}$\hss}}}
\begingroup\lccode`\~=`\- \lowercase{\endgroup\let~}\lycee@moins
\newcommand\lycee@moinsepais{\mathcode`\-="8000 }
\AtBeginDocument{\mathchardef\lycee@moinsorig=\mathcode`\-\relax}
\AtBeginDocument{\catcode`\_=12 \mathcode`\_="8000
  \catcode`\^=12 \mathcode`\^="8000 }
\newcommand\lycee@exposants{%
  \fontdimen13\textfont2=.48\fontdimen6\textfont2
  \fontdimen14\textfont2=.48\fontdimen6\textfont2
  \fontdimen15\textfont2=.48\fontdimen6\textfont2
  \fontdimen13\scriptfont2=.48\fontdimen6\scriptfont2
  \fontdimen14\scriptfont2=.48\fontdimen6\scriptfont2
  \fontdimen15\scriptfont2=.48\fontdimen6\scriptfont2}
\every@math@size\expandafter{\the\every@math@size\lycee@exposants}
\newlength{\retraitcalculs}
\setlength{\retraitcalculs}{2em}

\RequirePackage{tikz}
\usetikzlibrary{arrows.meta,calc,babel,shapes.misc}
\RequirePackage[most]{tcolorbox}
\tcbuselibrary{skins,breakable}

\definecolor{coulDef}{HTML}{1F4E79}
\definecolor{coulProp}{HTML}{7030A0}
\definecolor{coulRem}{HTML}{C55A11}
\definecolor{coulEx}{HTML}{2E7D32}
\definecolor{coulExo}{HTML}{B03A2E}
\definecolor{coulPart}{HTML}{1F4E79}
\definecolor{coulMeth}{HTML}{1B6E4F}
\definecolor{coulCourbe}{HTML}{0B5ED7}
\definecolor{coulCourbeB}{HTML}{C00000}
\definecolor{coulCourbeC}{HTML}{2E7D32}
\definecolor{coulGrille}{HTML}{8C8C8C}
\definecolor{coulRep}{HTML}{1B6E4F}
\definecolor{coulBar}{HTML}{2E7D32}

\newcommand{\R}{\mathbb{R}}
\newcommand{\Rs}{\mathbb{R}^{*}}
\renewcommand{\le}{\leqslant}
\renewcommand{\ge}{\geqslant}
\newcommand{\vect}[1]{\overrightarrow{#1}}
\newcommand{\norme}[1]{\left\lVert #1 \right\rVert}
\newcommand{\intervalle}[2]{\left[#1\,;#2\right]}
\RequirePackage{cancel}
\renewcommand{\CancelColor}{\color{coulExo}}
\newcommand{\fois}[1]{\times{\color{coulExo}#1}}
\newcommand{\barre}[1]{\cancel{#1}}

\tikzset{grille/.style={coulGrille,line width=0.4pt}}
\newcommand{\quadrillage}[4]{\draw[grille,step=1] (#1,#2) grid (#3,#4);}
\tikzset{
  croix/.style={cross out,draw,line width=0.6pt,inner sep=0pt,minimum size=4pt},
  croixdroite/.style={croix,rotate=45,line width=0.8pt},
}
\newcommand{\pt}[3]{\path (#1) node[croix]{} node[#2,font=\small,inner sep=2.5pt]{$#3$};}

\newcommand{\lycee@repere}[4]{%
  \draw[grille,step=1] (#1,#2) grid (#3,#4);
  \draw[-{Stealth[length=2mm]},thick] (#1,0)--({#3+0.4},0) node[below left=-1pt]{$x$};
  \draw[-{Stealth[length=2mm]},thick] (0,#2)--(0,{#4+0.4}) node[below left=-1pt]{$y$};
  \node[below left,font=\scriptsize,inner sep=1.5pt] at (0,0) {$0$};}
\newcommand{\lycee@gradx}[1]{\foreach \k/\l in {#1}{%
  \draw (\k,0.07)--(\k,-0.07) node[below,font=\scriptsize,inner sep=1.5pt]{$\l$};}}
\newcommand{\lycee@grady}[1]{\foreach \k/\l in {#1}{%
  \draw (0.07,\k)--(-0.07,\k) node[left,font=\scriptsize,inner sep=1.5pt]{$\l$};}}
\newcommand{\lycee@intff}[2]{\left[#1\,;#2\right]}
\newcommand{\lycee@intfo}[2]{\left[#1\,;#2\right[}
\newcommand{\lycee@intof}[2]{\left]#1\,;#2\right]}
\newcommand{\lycee@intoo}[2]{\left]#1\,;#2\right[}
\newcommand{\lycee@ens}[1]{\left\{#1\right\}}
\AtBeginDocument{%
  \providecommand{\repere}{\lycee@repere}%
  \providecommand{\gradx}{\lycee@gradx}%
  \providecommand{\grady}{\lycee@grady}%
  \providecommand{\Cf}{\mathcal{C}\sb{\scriptscriptstyle f}}%
  \providecommand{\Cg}{\mathcal{C}\sb{\scriptscriptstyle g}}%
  \providecommand{\intff}{\lycee@intff}%
  \providecommand{\intfo}{\lycee@intfo}%
  \providecommand{\intof}{\lycee@intof}%
  \providecommand{\intoo}{\lycee@intoo}%
  \providecommand{\ens}{\lycee@ens}}

\newlist{questions}{enumerate}{3}
\setlist[questions,1]{label=\arabic*\textdegree),leftmargin=*,itemsep=2pt,topsep=3pt}
\setlist[questions,2]{label=\alph*),leftmargin=*,itemsep=1pt,topsep=2pt}
\setlist[questions,3]{label=\roman*),leftmargin=*,itemsep=1pt,topsep=1pt}
\setlist[itemize]{label=\textbullet,leftmargin=*,itemsep=2pt,topsep=3pt}

\newcommand{\besoinplace}[1]{\par\penalty\@M\begingroup
  \dimen@=#1\relax\dimen@ii=\pagegoal\advance\dimen@ii-\pagetotal
  \ifdim\dimen@>\dimen@ii\ifdim\dimen@ii>\z@\vfil\fi\break\fi\endgroup}

\newcommand{\lycee@pied}{\thepage/\pageref{LastPage}}
\newcommand{\entete}[2]{%
  \pagestyle{fancy}\fancyhf{}%
  \fancyhead[L]{\footnotesize\color{coulPart}\textbf{#1}}%
  \fancyhead[R]{\footnotesize\color{coulPart}\textbf{#2}}%
  \fancyfoot[C]{\footnotesize\lycee@pied}%
  \renewcommand{\headrulewidth}{0.4pt}%
  \renewcommand{\headrule}{\hbox to\headwidth{%
    \color{coulPart!40}\leaders\hrule height \headrulewidth\hfill}}}

\RequirePackage{ccicons}
\newcommand{\lycee@licence}{{\scriptsize\color{gray}%
  \copyright\ Luc Destombe --- \ccbyncsa\ CC BY-NC-SA 4.0}}
\newif\iflycee@licence \lycee@licencetrue
\newcommand{\sanslicence}{\lycee@licencefalse}
\AtBeginDocument{\iflycee@licence\AddToHookNext{shipout/foreground}{%
  \put(0,0){\rlap{\kern\dimexpr1in+\hoffset+\oddsidemargin+\textwidth\relax
    \llap{\raisebox{-\dimexpr1in+\voffset+\topmargin+\headheight+\headsep
      +\textheight+\footskip\relax}{\lycee@licence}}}}}\fi}

\newcommand{\formulesserrees}{%
  \setlength{\abovedisplayskip}{3pt}\setlength{\belowdisplayskip}{3pt}%
  \setlength{\abovedisplayshortskip}{2pt}\setlength{\belowdisplayshortskip}{3pt}}

\setlength{\parindent}{0pt}

\RequirePackage[implicit=false,hidelinks,pdfencoding=auto,psdextra,bookmarksopen,
  bookmarksopenlevel=1,pdfcreator={},pdfproducer={}]{hyperref}
\RequirePackage{bookmark}
\hypersetup{pdfauthor={Luc Destombe},pdfkeywords={CC BY-NC-SA 4.0}}
\newcounter{lycee@signet}
\newcommand{\signet}[3][]{\stepcounter{lycee@signet}%
  \edef\lycee@dest{\if\relax\detokenize{#1}\relax
    signet.\arabic{lycee@signet}\else#1\fi}%
  \hypertarget{\lycee@dest}{}\bookmark[level=#2,dest=\lycee@dest]{#3}}
\pdfstringdefDisableCommands{%
  \def\R{R}\def\vect#1{#1}\def\Cf{Cf}\def\Cg{Cg}\def\fois#1{×#1}%
  \def\barre#1{#1}\def\intervalle#1#2{[#1;#2]}\def\intff#1#2{[#1;#2]}%
  \def\textsuperscript#1{#1}\def\dfrac#1#2{#1/#2}\def\frac#1#2{#1/#2}%
  \def\vec#1{#1⃗}}%
\RequirePackage[official]{eurosym}

\tcbset{exercice corrige/.style={
  enhanced,breakable,before upper=\raggedright,
  colback=white,colframe=coulExo,
  boxrule=0.8pt,arc=2pt,
  left=6pt,right=6pt,top=8pt,bottom=5pt,
  before skip=9pt,after skip=4pt,
  coltitle=white,fonttitle=\bfseries\small,
  attach boxed title to top left={xshift=8pt,yshift=-\tcboxedtitleheight/2},
  boxed title style={colback=coulExo,arc=2pt,boxrule=0pt}}}
\newcounter{exo}
\newtcolorbox{exercice}[1][]{exercice corrige,
  phantom={\signet[exercice-\arabic{exo}]{2}{Exercice \arabic{exo}}},
  title={Exercice \theexo\ifx\relax#1\relax\else{} --- #1\fi}}
\newenvironment{exo}[1][\relax]{\refstepcounter{exo}\begin{exercice}[#1]}{\end{exercice}}

\RequirePackage{cuted}
\geometry{top=1.6cm,bottom=1cm,left=1cm,right=1cm,headheight=14pt,footskip=0.6cm}

\newcommand{\entetecorrige}[3]{%
  \begin{strip}
  \begin{center}
    \begin{tcolorbox}[enhanced,width=0.92\linewidth,colback=white,colframe=coulPart,
        boxrule=1.6pt,arc=3pt,halign=center,top=5pt,bottom=5pt]
      {\LARGE\bfseries\color{coulPart} #1}\\[3pt]
      {\normalsize #2}
    \end{tcolorbox}
  \end{center}
  \if\relax\detokenize{#3}\relax\else

  \vspace{2pt}
  \noindent
  \begin{tcolorbox}[enhanced,colback=white,colframe=black!35,boxrule=0.5pt,arc=2pt,
      left=7pt,right=7pt,top=4pt,bottom=4pt]
  {\small #3}
  \end{tcolorbox}
  \vspace{6pt}
  \fi
  \end{strip}}

\newcounter{partie}
\newcommand{\partiebandeau}[1]{%
  \besoinplace{8\baselineskip}\vspace{6pt}%
  \begin{tcolorbox}[enhanced,colback=coulPart!8,colframe=coulPart,boxrule=0pt,
      leftrule=3pt,arc=0pt,left=5pt,right=4pt,top=3pt,bottom=3pt,
      before skip=4pt,after skip=2pt]
    \bfseries\color{coulPart}#1\signet{1}{#1}
  \end{tcolorbox}}
\newcommand{\partie}[1]{\stepcounter{partie}\partiebandeau{\Roman{partie}) #1}}
\newcommand{\partiesansnum}[1]{\partiebandeau{#1}}

\newcommand{\res}[1]{%
  \tcbox[on line,colback=coulPart!7,colframe=coulPart,boxrule=0.5pt,arc=1pt,
    boxsep=0pt,left=3pt,right=3pt,top=1pt,bottom=2pt]{$#1$}}
\newcommand{\calc}[1]{\par\smallskip\noindent$\displaystyle #1$\par}
\newenvironment{calculs}{\par\smallskip\noindent$\displaystyle\begin{aligned}[t]}%
  {\end{aligned}$\par\smallskip}
\newcommand{\rem}[1]{\par\smallskip{\small\itshape\color{coulPart}#1}\par}
\newcommand{\deux}[2]{\par\noindent\makebox[0.5\linewidth][l]{#1}#2}

\setlist[questions,1]{label=\arabic*\textdegree),leftmargin=*,itemsep=2pt,topsep=2pt}
\setlist[questions,2]{label=\alph*),leftmargin=*,itemsep=2pt,topsep=1pt}
\newlist{reponses}{enumerate}{1}
\setlist[reponses]{label=\alph*),leftmargin=*,itemsep=2pt,topsep=2pt}

\setlength{\parskip}{1pt}
\setlength{\columnsep}{0.6cm}
\raggedbottom
\formulesserrees
\AtBeginDocument{\formulesserrees}

\makeatother
\entete{Chapitre 1 --- Fractions et puissances --- Corrigé des exercices}{Seconde}

\let\deuxpaquet\deux
\renewcommand{\deux}[2]{\deuxpaquet{#1}{#2}\par}
\clubpenalty=10000
\widowpenalty=10000

\begin{document}

\entetecorrige{CORRIGÉ --- FRACTIONS ET PUISSANCES}
  {Chapitre 1 : fiche d'exercices --- Seconde}
  {Les résultats sont encadrés ; les remarques en bleu signalent les erreurs
   fréquentes et les méthodes à retenir. Ce qui n'est pas exigé est en gris.}

\partie{Nombres relatifs et priorités}

\begin{exo}[Calculs sur les relatifs]
\deux{$A=\res{-16}$}{$B=14+11=\res{25}$}
\deux{$C=\res{-42}$}{$D=\res{40}$}
\deux{$E=\dfrac{27}{-9}=\res{-3}$}{$F=-13-4=\res{-17}$}
\rem{Soustraire $-11$, c'est ajouter $11$ ; deux facteurs négatifs donnent un produit
positif.}
\end{exo}

\begin{exo}[Priorités de calcul]
\begin{calculs}
A&=64-(-18)\\
A&=\res{82}\\
B&=-7-16\\
B&=\res{-23}\\
C&=6+(-4)\times 4\\
C&=6-16=\res{-10}\\
D&=-36\div(-4)\times 4\\
D&=9\times 4=\res{36}
\end{calculs}
\rem{Dans $B$, la puissance passe avant la soustraction. Dans $D$, division et produit
se font de gauche à droite.}
\end{exo}

\begin{exo}[Parenthèses]
\begin{calculs}
A&=20-4\times(-5)\\
A&=20+20=\res{40}\\
B&=-15+8=\res{-7}\\
C&=-48\div(-6)=\res{8}\\
D&=3-2\times(5+5)\\
D&=3-20=\res{-17}
\end{calculs}
\rem{Dans $D$, on commence par la parenthèse la plus intérieure : $1-6=-5$, puis
$5-(-5)=10$.}
\end{exo}

\partie{Additionner et soustraire des fractions}

\begin{exo}[Même dénominateur]
\deux{$A=\res{\dfrac75}$}{$B=\dfrac{-2+7}{9}=\res{\dfrac59}$}
\deux{$C=\dfrac{5-9}{7}=\res{-\dfrac47}$}{$D=\dfrac{-3+8}{11}=\res{\dfrac{5}{11}}$}
\deux{$E=\dfrac{6}{12}=\res{\dfrac12}$}{$F=\dfrac{-12}{4}=\res{-3}$}
\rem{Le dénominateur ne change pas : on n'additionne jamais les dénominateurs.}
\end{exo}

\begin{exo}[Dénominateurs différents]
\begin{calculs}
A&=\dfrac{8}{20}+\dfrac{15}{20}=\res{\dfrac{23}{20}}\\
B&=\dfrac{10}{12}+\dfrac{3}{12}=\res{\dfrac{13}{12}}\\
C&=\dfrac{9}{12}+\dfrac{2}{12}=\res{\dfrac{11}{12}}\\
D&=\dfrac{9}{24}-\dfrac{20}{24}=\res{-\dfrac{11}{24}}\\
E&=-\dfrac{10}{15}-\dfrac{12}{15}=\res{-\dfrac{22}{15}}\\
F&=-\dfrac{14}{12}+\dfrac{15}{12}=\res{\dfrac{1}{12}}
\end{calculs}
\rem{Pour $B$, $C$, $D$, $F$, le plus petit dénominateur commun n'est pas le produit
des dénominateurs : $12$ est un multiple de $6$ et de $4$.}
\end{exo}

\begin{exo}[Avec un nombre entier]
\deux{$A=\dfrac{12}{4}-\dfrac{5}{4}=\res{\dfrac74}$}{$B=-\dfrac{10}{5}-\dfrac{3}{5}=\res{-\dfrac{13}{5}}$}
\deux{$C=-\dfrac{24}{6}+\dfrac{17}{6}=\res{-\dfrac76}$}{$D=\dfrac{9}{2}-\dfrac{8}{2}=\res{\dfrac12}$}
\end{exo}

\begin{exo}[Avec une lettre]
\begin{calculs}
A&=\dfrac{12x}{4}-\dfrac{x}{4}=\res{\dfrac{11x}{4}}\\
B&=\dfrac{2x}{5}-\dfrac{15}{5}=\res{\dfrac{2x-15}{5}}\\
C&=\dfrac{6}{3}-\dfrac{x-1}{3}\\
C&=\dfrac{6-(x-1)}{3}=\res{\dfrac{7-x}{3}}\\
D&=\dfrac{2x+1-3x}{3}=\res{\dfrac{1-x}{3}}
\end{calculs}
\rem{Dans $C$, le moins porte sur tout le numérateur : $-(x-1)=-x+1$.}
\end{exo}

\partie{Multiplier et diviser des fractions}

\begin{exo}[Simplifier avant de calculer]
\begin{calculs}
A&=-\dfrac{2\times 2\times 3\times 3}{2\times 3\times 2\times 5}=\res{-\dfrac35}\\
B&=\dfrac{2\times 3\times 4\times 5\times 7}{3\times 7\times 2\times 5}=\res{4}\\
C&=-\dfrac{2\times 13\times 3\times 3}{3\times 5\times 3\times 13}=\res{-\dfrac25}\\
D&=\dfrac{1\times 2\times 7\times 3\times 5}{2\times 5\times 3\times 7}=\res{1}
\end{calculs}
\rem{On fixe d'abord le signe : deux signes moins dans $B$ et $D$ (positif), un seul
dans $A$ et $C$ (négatif). Dans $D$, tout se barre : il reste $1$, et non $0$.}
\end{exo}

\begin{exo}[Simplifier avec des lettres]
\begin{calculs}
A&=\dfrac{3\times 5\times 2\times 4\times 2\times 7}{5\times 7\times 2\times 3\times 4}=\res{2}\\
B&=\dfrac{3\times 4\times 5\times x}{4\times 5}=\res{3x}\\
C&=\dfrac{2\times 3\times 5\times a\times b}{2\times 5\times b\times 3\times 3}=\res{\dfrac{a}{3}}\\
D&=\dfrac{2\times 2\times a\times 3\times 3}{2\times 3\times 2\times a}=\res{3}
\end{calculs}
\end{exo}

\begin{exo}[Produits]
\deux{$A=\res{-\dfrac{15}{28}}$}{$B=\res{-\dfrac{8}{15}}$}
\deux{$C=\res{-\dfrac{15}{8}}$}{$D=\dfrac73\times\dfrac{-2}{5}=\res{-\dfrac{14}{15}}$}
\deux{$E=\res{\dfrac{16}{45}}$}{$F=\res{-\dfrac{15}{56}}$}
\rem{$F$ compte trois signes moins : le produit est négatif.}
\end{exo}

\begin{exo}[Quotients]
\begin{calculs}
A&=\dfrac25\times\dfrac73=\res{\dfrac{14}{15}}\\
B&=\dfrac54\times\dfrac37=\res{\dfrac{15}{28}}\\
C&=\dfrac34\times\dfrac16=\res{\dfrac18}\\
D&=6\times\dfrac{5}{-2}=\res{-15}\\
E&=\dfrac49\times\dfrac{2}{-5}=\res{-\dfrac{8}{45}}\\
F&=3\times\dfrac{7}{-4}=\res{-\dfrac{21}{4}}
\end{calculs}
\rem{Diviser par $6$, c'est multiplier par $\frac16$.}
\end{exo}

\partie{Fractions et priorités de calcul}

\begin{exo}[Respecter les priorités]
\begin{calculs}
A&=\dfrac{5}{12}+\dfrac{9}{12}=\dfrac{14}{12}\\
A&=\res{\dfrac76}\\
B&=\dfrac{7}{10}+\dfrac{4}{10}=\res{\dfrac{11}{10}}\\
C&=-\dfrac{10}{12}+\dfrac{9}{12}-\dfrac{8}{12}\\
C&=-\dfrac{9}{12}=\res{-\dfrac34}\\
D&=\dfrac{9}{14}+\dfrac{1}{2}\\
D&=\dfrac{9}{14}+\dfrac{7}{14}=\res{\dfrac87}
\end{calculs}
\rem{Dans $D$, le produit d'abord : $\frac57\times\frac{7}{10}=\frac{5}{10}=\frac12$.}
\end{exo}

\begin{exo}[Parenthèses et lettres]
\begin{calculs}
A&=\dfrac45\times\dfrac14=\res{\dfrac15}\\
B&=\dfrac{3}{6}\times\dfrac{3}{4}\\
B&=\dfrac12\times\dfrac34=\res{\dfrac38}\\
C&=\dfrac{9x}{12}-\dfrac{2(x-2)}{12}\\
C&=\dfrac{9x-2x+4}{12}=\res{\dfrac{7x+4}{12}}\\
D&=\dfrac{3\times 8\times x\times x}{4\times 9}=\res{\dfrac{2x^{2}}{3}}
\end{calculs}
\rem{Dans $C$ : $-2\times(-2)=+4$.}
\end{exo}

\begin{exo}[Fractions de fractions]
\begin{calculs}
A&=\dfrac34\times\dfrac65=\dfrac{18}{20}\\
A&=\res{\dfrac{9}{10}}\\
B&=\dfrac{6}{35}\times\dfrac{14}{9}=\dfrac{2\times 3\times 2\times 7}{5\times 7\times 3\times 3}\\
B&=\res{\dfrac{4}{15}}\\
C&=\dfrac74\times\dfrac{1}{-14}=\res{-\dfrac18}
\end{calculs}
Pour $D$, numérateur et dénominateur à part :
\begin{calculs}
2-\dfrac35&=\dfrac{10}{5}-\dfrac35=\dfrac75\\
\dfrac34-\dfrac25&=\dfrac{15}{20}-\dfrac{8}{20}=\dfrac{7}{20}\\
D&=\dfrac75\times\dfrac{20}{7}=\res{4}
\end{calculs}
\end{exo}

\partie{Produit en croix}

\begin{exo}[Justifier le produit en croix]
1) Les deux fractions sont égales, et leurs dénominateurs $4$ et $2$ ne sont pas nuls :
d'après le produit en croix, $2\times(3x-1)=4\times(x+3)$.

\smallskip
2)
\begin{calculs}
6x-2&=4x+12\\
6x-4x&=12+2\\
2x&=14\\
x&=\res{7}
\end{calculs}
\rem{Vérification : $\frac{20}{4}=5$ et $\frac{10}{2}=5$.}
\end{exo}

\begin{exo}[Premières équations]
\textbf{a)}
\begin{calculs}
3x&=20\\
x&=\res{\dfrac{20}{3}}
\end{calculs}
\textbf{b)}
\begin{calculs}
5(4-x)&=3(x+2)\\
20-5x&=3x+6\\
14&=8x\\
x&=\dfrac{14}{8}=\res{\dfrac74}
\end{calculs}
\textbf{c)}
\begin{calculs}
6(2+x)&=-3\times 5\\
12+6x&=-15\\
6x&=-27\\
x&=-\dfrac{27}{6}=\res{-\dfrac92}
\end{calculs}
\textbf{d)}
\begin{calculs}
-2(2x-1)&=3(x+4)\\
-4x+2&=3x+12\\
-10&=7x\\
x&=\res{-\dfrac{10}{7}}
\end{calculs}
\rem{Un dénominateur négatif garde son signe dans le produit en croix.}
\end{exo}

\begin{exo}[S'entraîner]
\textbf{a)} $8x=15$, donc $x=\res{\dfrac{15}{8}}$.\par
\textbf{b)} $9x=24$, donc $x=\dfrac{24}{9}=\res{\dfrac83}$.\par
\textbf{c)}
\begin{calculs}
4(x+2)&=5(2x-3)\\
4x+8&=10x-15\\
23&=6x\\
x&=\res{\dfrac{23}{6}}
\end{calculs}
\textbf{d)}
\begin{calculs}
4(2x+1)&=3(3x-2)\\
8x+4&=9x-6\\
x&=\res{10}
\end{calculs}
\rem{Vérification de d) : $\frac{21}{3}=7$ et $\frac{28}{4}=7$.}
\end{exo}

\begin{exo}[Plusieurs fractions (1)]
\textbf{a)} ($\times 5$)
\begin{calculs}
5x+x&=90\\
x&=\res{15}
\end{calculs}
\textbf{b)} ($\times 12$)
\begin{calculs}
4x-3x&=60\\
x&=\res{60}
\end{calculs}
\textbf{c)} ($\times 15$)
\begin{calculs}
5t-3(t+1)&=15\\
2t-3&=15\\
t&=\res{9}
\end{calculs}
\textbf{d)} ($\times 6$)
\begin{calculs}
3(3t-1)+2(t-2)&=12t\\
11t-7&=12t\\
t&=\res{-7}
\end{calculs}
\rem{Chaque terme est multiplié, y compris celui sans fraction ($x$ dans a), $1$ dans
c), $2t$ dans d)).}
\end{exo}

\begin{exo}[Plusieurs fractions (2)]
\textbf{a)} ($\times 12$)
\begin{calculs}
6y+9y&=4y-264\\
11y&=-264\\
y&=\res{-24}
\end{calculs}
\textbf{b)} ($\times 12$)
\begin{calculs}
4(y+7)&=6(y-1)-3y\\
4y+28&=3y-6\\
y&=\res{-34}
\end{calculs}
\textbf{c)} ($\times 36$)
\begin{calculs}
12(z-4)-3(z-6)&=4(2z-15)\\
9z-30&=8z-60\\
z&=\res{-30}
\end{calculs}
\textbf{d)} ($\times 10$)
\begin{calculs}
2(3z+2)-5(2z-1)&=3-6z\\
-4z+9&=3-6z\\
2z&=-6\\
z&=\res{-3}
\end{calculs}
\rem{Attention aux signes devant les parenthèses : $-3\times(-6)=+18$,
$-5\times(-1)=+5$.}
\end{exo}

\partie{Fractions et problèmes}

\begin{exo}[Fraction d'une quantité]
\deux{a) $\dfrac34\times 120=\res{90\text{ €}}$}{b) $\dfrac25\times 35=\res{14\text{ kg}}$}
\deux{c) $\dfrac56\times 90=\res{75\text{ min}}$}{d) $\dfrac23\times\dfrac34\times 48=\res{24}$}
\rem{En d), on peut calculer $\frac23\times\frac34=\frac12$ d'abord.}
\end{exo}

\begin{exo}[Le festival]
\begin{calculs}
\text{prévente}&=\dfrac56\times 1~200=1~000\\
\text{restantes}&=1~200-1~000=200\\
\text{sur place}&=\dfrac34\times 200=150\\
\text{invendues}&=200-150=\res{50}
\end{calculs}
\rem{Les $\frac34$ portent sur les places restantes, pas sur les $1~200$.}
\end{exo}

\begin{exo}[L'anniversaire de Lina]
a)
\begin{calculs}
\text{reste}&=1-\dfrac13=\dfrac23\\
\text{vêtements}&=\dfrac38\times\dfrac23=\dfrac14\\
\text{reste final}&=\dfrac23-\dfrac14=\res{\dfrac{5}{12}}
\end{calculs}
b) Jeux vidéo : $\frac13\times 120=\res{40\text{ €}}$ ; vêtements :
$\frac14\times 120=\res{30\text{ €}}$ ; il reste $\frac{5}{12}\times 120=\res{50\text{ €}}$.
\rem{Vérification : $40+30+50=120$.}
\end{exo}

\begin{exo}[Le budget d'une association]
a)
\begin{calculs}
\text{reste}&=1-\dfrac25=\dfrac35\\
\text{sorties}&=\dfrac34\times\dfrac35=\res{\dfrac{9}{20}}\\
\text{frais}&=\dfrac35-\dfrac{9}{20}=\dfrac{12}{20}-\dfrac{9}{20}\\
\text{frais}&=\res{\dfrac{3}{20}}
\end{calculs}
b) Les $\frac{3}{20}$ du budget font $270$~€, donc $\frac{1}{20}$ fait $90$~€ et le
budget est de $20\times 90=\res{1~800\text{ €}}$.
\rem{Vérification : matériel $720$~€, sorties $810$~€, frais $270$~€ ; total $1~800$~€.}
\end{exo}

\partie{Puissances}

\begin{exo}[Calculer des puissances]
\deux{$A=\res{81}$}{$B=\res{16}$}
\deux{$C=\dfrac{1}{4^{2}}=\res{\dfrac{1}{16}}$}{$D=\dfrac{1}{(-3)^{3}}=\res{-\dfrac{1}{27}}$}
\deux{$E=\res{0{,}001}$}{$F=\res{1}$}
\deux{$G=\res{\dfrac17}$}{$H=\res{-9}$}
\rem{$4^{-2}$ n'est pas négatif : c'est l'inverse de $4^{2}$.}
\end{exo}

\begin{exo}[Écriture décimale]
\deux{$A=\res{0{,}001}$}{$B=\res{0{,}000~000~1}$}
\deux{$C=\res{100~000}$}{$D=\res{1}$}
\end{exo}

\begin{exo}[Puissances et priorités]
\deux{$A=2+25=\res{27}$}{$B=7^{2}=\res{49}$}
\deux{$C=2-25=\res{-23}$}{$D=27+16=\res{43}$}
\deux{$E=16\times 9=\res{144}$}{$F=-27+32=\res{5}$}
\rem{$A$ et $B$ : la parenthèse change tout.}
\end{exo}

\begin{exo}[Puissances et priorités (suite)]
\begin{calculs}
A&=(27-16)\times 100=\res{1~100}\\
B&=9+16\times 0{,}1=\res{10{,}6}\\
C&=16-100\div 25=\res{12}\\
D&=100\times(-1)-(-32)\\
D&=-100+32=\res{-68}
\end{calculs}
\end{exo}

\begin{exo}[Des calculs plus longs]
\begin{calculs}
A&=20-9\times(-2)^{2}\\
A&=20-36=\res{-16}\\
B&=\dfrac{125-80}{81-54}=\dfrac{45}{27}\\
B&=\res{\dfrac53}\\
C&=\dfrac{42-32}{40-36}=\dfrac{10}{4}\\
C&=\res{\dfrac52}
\end{calculs}
\end{exo}

\begin{exo}[Des résultats en fraction]
\begin{calculs}
A&=\dfrac19+\dfrac{2}{10}=\dfrac{5}{45}+\dfrac{9}{45}\\
A&=\res{\dfrac{14}{45}}\\
B&=\dfrac{9+8}{\frac{1}{10}}=17\times 10=\res{170}\\
C&=\dfrac14-\dfrac{1}{10}=\dfrac{5}{20}-\dfrac{2}{20}\\
C&=\res{\dfrac{3}{20}}\\
D&=16\times\dfrac{1}{16}=\res{1}
\end{calculs}
\end{exo}

\partie{Règles de calcul sur les puissances}

\begin{exo}[Produits]
\deux{$A=3^{4+7}=\res{3^{11}}$}{$B=\res{(-5)^{8}}$}
\deux{$C=7^{2-6}=\res{7^{-4}}$}{$D=x^{3+1+5}=\res{x^{9}}$}
\rem{$x=x^{1}$ : on n'oublie pas cet exposant.}
\end{exo}

\begin{exo}[Quotients]
\deux{$A=\res{6^{5}}$}{$B=\res{(-3)^{3}}$}
\deux{$C=\res{y^{3}}$}{$D=2^{-4-3}=\res{2^{-7}}$}
\deux{$E=4^{-1-6}=\res{4^{-7}}$}{$F=x^{-5-(-9)}=\res{x^{4}}$}
\end{exo}

\begin{exo}[Produits et quotients]
\deux{$A=\dfrac{3^{14}}{3^{11}}=\res{3^{3}}$}{$B=\dfrac{x^{5}}{x^{11}}=\res{x^{-6}}$}
\deux{$C=\dfrac{2^{11}}{2^{1}}=\res{2^{10}}$}{$D=\dfrac{y^{10}}{y^{7}}=\res{y^{3}}$}
\end{exo}

\begin{exo}[Toutes les règles]
\begin{calculs}
A&=\dfrac{5^{-4}\times 3^{1}}{5^{-3}\times 3^{-3}}\\
A&=5^{-4-(-3)}\times 3^{1-(-3)}\\
A&=5^{-1}\times 3^{4}=\res{\dfrac{81}{5}}\\
B&=\dfrac{3^{0}\times 10^{7}}{3^{-3}\times 10^{-8}}\\
B&=3^{3}\times 10^{15}=\res{2{,}7\times 10^{16}}
\end{calculs}
\rem{$3^{4}\times 3^{-4}=3^{0}=1$ ; $(10^{2})^{-4}=10^{-8}$.}
\end{exo}

\begin{exo}[Calculer et simplifier]
\begin{calculs}
A&=\dfrac{2^{12}\times 10^{-7}}{2^{10}\times 10^{-5}}\\
A&=2^{2}\times 10^{-2}=\res{4\times 10^{-2}}\\
B&=\dfrac{b^{2}\times b^{-9}}{b^{8}}=\dfrac{b^{-7}}{b^{8}}\\
B&=\res{b^{-15}}
\end{calculs}
\end{exo}

\partie{Notation scientifique}

\begin{exo}[Entiers et décimaux]
\deux{$A=\res{2{,}74\times 10^{5}}$}{$B=\res{6{,}1\times 10^{-4}}$}
\deux{$C=\res{3{,}82\times 10^{7}}$}{$D=\res{2{,}9\times 10^{-7}}$}
\deux{$E=\res{5{,}63\times 10^{4}}$}{$F=\res{7{,}308\times 10^{-5}}$}
\end{exo}

\begin{exo}[Produits par une puissance de 10]
\begin{calculs}
A&=5{,}32\times 10^{2}\times 10^{3}=\res{5{,}32\times 10^{5}}\\
B&=6{,}4\times 10^{1}\times 10^{-5}=\res{6{,}4\times 10^{-4}}\\
C&=4{,}817\times 10^{2}\times 10^{8}=\res{4{,}817\times 10^{10}}\\
D&=3{,}6\times 10^{-2}\times 10^{-4}=\res{3{,}6\times 10^{-6}}
\end{calculs}
\end{exo}

\begin{exo}[Un calcul jusqu'au bout]
\begin{calculs}
E&=\dfrac{35\times 6}{14}\times\dfrac{10^{-2}\times 10^{9}}{10^{5}}\\
E&=15\times 10^{2}\\
E&=\res{1{,}5\times 10^{3}}
\end{calculs}
\rem{$15\times 10^{2}$ n'est pas en notation scientifique : un seul chiffre avant la
virgule.}
\end{exo}

\begin{exo}[Deux expressions]
1)
\begin{calculs}
A&=\dfrac{3}{12}\times 10^{1500-1502}\\
A&=0{,}25\times 10^{-2}=\res{2{,}5\times 10^{-3}}
\end{calculs}
2)
\begin{calculs}
B&=\dfrac{640-34}{101}=\dfrac{606}{101}\\
B&=\res{6}
\end{calculs}
$B$ est bien un nombre entier.
\end{exo}

\begin{exo}[La lumière de la Lune]
1) $384~000\text{ km}=384~000~000\text{ m}=\res{3{,}84\times 10^{8}\text{ m}}$.

\smallskip
2)
\begin{calculs}
t&=\dfrac{3{,}84\times 10^{8}}{3\times 10^{8}}\\
t&=\res{1{,}28\text{ s}}
\end{calculs}
\rem{Le temps est la distance divisée par la distance parcourue en une seconde.}
\end{exo}

\end{document}
